\documentclass[11pt,a4paper]{article}

\usepackage[utf8]{inputenc}
\usepackage[T1]{fontenc}
\usepackage{lmodern}
\usepackage[margin=0.76in]{geometry}
\usepackage[table]{xcolor}
\usepackage{array}
\usepackage{tabularx}
\usepackage{booktabs}
\usepackage{enumitem}
\usepackage{titlesec}
\usepackage{fancyhdr}
\usepackage{lastpage}
\usepackage{hyperref}

\definecolor{TextInk}{HTML}{233142}
\definecolor{HeadingBlue}{HTML}{1D3557}
\definecolor{AccentTeal}{HTML}{2A7F92}
\definecolor{AccentGold}{HTML}{C58B4E}
\definecolor{TableStripe}{HTML}{F3F8FA}
\definecolor{SoftBlue}{HTML}{EEF6F8}
\definecolor{SoftGreen}{HTML}{EEF7F0}
\definecolor{SoftRed}{HTML}{FBF0EF}

\hypersetup{
  colorlinks=true,
  linkcolor=HeadingBlue,
  urlcolor=AccentTeal,
  pdftitle={IB Geography SL Paper 1 Simulation Set 3 Marking Guide - 2026-05-07},
  pdfauthor={IB Geography SL Preparation}
}

\color{TextInk}
\setlength{\parskip}{0.46em}
\setlength{\parindent}{0pt}
\setlength{\tabcolsep}{5pt}
\renewcommand{\arraystretch}{1.18}
\setlist[itemize]{leftmargin=1.32em,itemsep=3pt,topsep=4pt}

\newcolumntype{Y}{>{\raggedright\arraybackslash}X}
\newcolumntype{L}[1]{>{\raggedright\arraybackslash}p{#1}}

\titleformat{\section}
  {\normalfont\huge\sffamily\bfseries\color{HeadingBlue}}
  {}{0pt}{}
  [\vspace{0.12em}{\color{AccentGold}\titlerule[1pt]}]
\titleformat{\subsection}
  {\normalfont\Large\sffamily\bfseries\color{AccentTeal}}
  {}{0pt}{}

\pagestyle{fancy}
\setlength{\headheight}{14pt}
\fancyhf{}
\fancyhead[L]{\sffamily\color{AccentTeal}Set 3 Marking Guide}
\fancyhead[R]{\sffamily\color{HeadingBlue}IB Geography SL}
\fancyfoot[C]{\sffamily\color{HeadingBlue}\thepage\ of \pageref{LastPage}}
\renewcommand{\headrulewidth}{0.4pt}

\newcommand{\term}[1]{\textcolor{HeadingBlue}{\textbf{#1}}}
\newcommand{\exam}[1]{\textcolor{AccentTeal}{\textbf{#1}}}
\newcommand{\boxnote}[3]{%
\noindent\colorbox{#1}{%
  \begin{minipage}{0.965\textwidth}
  \sffamily
  \textbf{\textcolor{HeadingBlue}{#2}} #3
  \end{minipage}
}}

\begin{document}

\begin{center}
  {\sffamily\bfseries\color{HeadingBlue}\LARGE Paper 1 Simulation Set 3 Marking Guide}\\[0.25em]
  {\sffamily\color{AccentTeal}\large Urban Environments + Food And Health}
\end{center}

\boxnote{SoftRed}{Use after the timed paper only.}{This guide is for
self-marking. It is not an official IB mark scheme. Award marks for accurate
geography, clear source use, valid explanation, named examples, and supported
judgment.}

\section{Strict Marking Rules}

\begin{itemize}
  \item Mark only what was written during the timed attempt.
  \item Source-based short answers must use exact categories, numbers, or
  phrases from the source.
  \item Explanation marks require a causal chain, not a list.
  \item Ten-mark answers must answer the command term and include a judgment.
\end{itemize}

\section{Urban Environments Mark Guide}

\subsection{Question 1: State Uneven Benefits [2]}

Award up to \term{2 marks}, one for each valid source-based point.

\begin{itemize}
  \item Property values increased by \term{64\%} in the waterfront zone but
  decreased by \term{4\%} in the vacant industrial fringe.
  \item The CBD has \term{89\%} of jobs within 30 minutes by transit, compared
  with only \term{23\%} for the vacant industrial fringe.
  \item The waterfront has only \term{14} affordable homes per 1000 residents,
  while the inner neighbourhood has \term{36}.
  \item The vacant industrial fringe still has derelict land, weak access, and
  limited private investment.
\end{itemize}

\subsection{Question 2: Describe Two Contrasts [4]}

Award up to \term{4 marks}: up to \term{2 marks} for each clear contrast using
both areas.

\rowcolors{2}{TableStripe}{white}
\begin{tabularx}{\textwidth}{L{0.24\textwidth}Y Y}
\toprule
\textbf{Contrast} & \textbf{Waterfront redevelopment zone} & \textbf{Vacant industrial fringe} \\
\midrule
Property values & Increased strongly by \term{64\%} & Decreased by
\term{4\%} \\
Jobs and access & \term{76\%} of jobs within 30 minutes by transit & Only
\term{23\%} within 30 minutes by transit \\
Housing and land use & New offices, leisure space, and high-income housing,
with \term{14} affordable homes per 1000 residents & Derelict land and limited
private investment, with \term{20} affordable homes per 1000 residents \\
Regeneration outcome & Visible reinvestment and higher-value functions &
Continued decline or weak regeneration \\
\bottomrule
\end{tabularx}
\rowcolors{2}{}{}

\subsection{Question 3: Explain Winners And Losers [4]}

Award up to \term{4 marks}: up to \term{2 marks} for each explained chain.

\begin{itemize}
  \item Regeneration can attract investment, jobs, and services to selected
  districts, so property owners and businesses may benefit more than renters.
  \item Rising land values and rents can displace lower-income residents or
  reduce affordability, creating social stress.
  \item Public and private investment may be spatially selective, leaving
  peripheral or industrial areas with weak access and derelict land.
  \item New jobs may not match local residents' skills, so economic gains can
  bypass the original community.
\end{itemize}

\subsection{Question 4(a): Gentrification Impacts [10]}

Credit named examples such as Detroit, London Docklands, New York, or a
class-specific case. Strong answers examine benefits such as reinvestment,
services, image, and tax base, and costs such as displacement, cultural change,
rent pressure, and uneven access.

\subsection{Question 4(b): Managing Decline Or Deprivation [10]}

Credit any valid named city. Detroit is a strong anchor if the answer evaluates
regeneration, affordable housing, job access, service improvement, reuse of
vacant land, and uneven outcomes. Curitiba is stronger for transport than for
post-industrial decline unless the answer is clearly framed around deprivation.

\section{Food And Health Mark Guide}

\subsection{Question 1: State Stakeholder Contrasts [2]}

Award up to \term{2 marks}, one for each valid source-based contrast.

\begin{itemize}
  \item TNC convenience-food expansion scores \term{5} for food availability
  but only \term{1} for diet quality.
  \item Government school meals have the highest equity reach, \term{5}, while
  TNC convenience-food expansion has only \term{3}.
  \item Smallholder co-operatives score \term{4} for sustainability, higher
  than TNC convenience-food expansion, at \term{2}.
  \item NGO WASH and vaccination programmes score only \term{1} for food
  availability but \term{4} for equity reach.
\end{itemize}

\subsection{Question 2: Describe Mixed Outcomes [4]}

Award up to \term{4 marks}: up to \term{2 marks} for each described mixed
outcome.

\begin{itemize}
  \item TNC convenience-food expansion may increase food availability, score
  \term{5}, but diet quality is weak, score \term{1}, because products may be
  ultra-processed.
  \item Government school meals score well for availability, diet quality, and
  equity, but the limitation is uneven funding and coverage.
  \item NGO WASH and vaccination programmes can improve health equity but do
  not directly increase food supply.
  \item Smallholder co-operatives have relatively strong diet quality and
  sustainability scores but may be limited by finance and scale.
\end{itemize}

\subsection{Question 3: Explain Reducing Inequality [4]}

Award up to \term{4 marks}: up to \term{2 marks} for each explained chain.

\begin{itemize}
  \item Governments can target vulnerable groups through school meals, public
  health campaigns, subsidies, and food standards, reducing unequal access to
  nutritious food.
  \item NGOs can improve WASH, vaccination, and health education, reducing
  disease exposure among poorer communities.
  \item Local co-operatives can support small farmers, improve local food
  availability, and keep more value in rural communities.
  \item Regulation of TNC marketing, sugar, salt, or food safety can reduce
  health risks linked to ultra-processed diets.
\end{itemize}

\subsection{Question 4(a): TNCs [10]}

Credit a named example such as Brazil and Nestle or another class-specific
TNC. Strong answers judge both sides: TNCs can improve distribution, jobs, and
availability, but may worsen diet quality, processed-food dependence, market
power, and health inequality.

\subsection{Question 4(b): Governments, NGOs, Or Local Communities [10]}

Credit named examples of school meals, WASH programmes, vaccination, small
farmer support, public health policy, or community food projects. Strong
answers evaluate cost, targeting, scale, sustainability, local participation,
and whether the strategy addresses causes rather than symptoms.

\newpage

\section{Generic 10-Mark Level Guide}

\rowcolors{2}{TableStripe}{white}
\begin{tabularx}{\textwidth}{L{0.14\textwidth}Y}
\toprule
\textbf{Marks} & \textbf{Descriptor} \\
\midrule
0 & No relevant geography, blank answer, or wholly incorrect response. \\
1--2 & Basic general statements. Little or no named example. Mainly descriptive
or asserted. \\
3--4 & Some relevant knowledge and a partly relevant example. Limited
explanation, balance, or evaluation. \\
5--6 & Clear, mostly accurate answer with a named example and some explanation.
Shows some balance, but evidence or judgment is uneven. \\
7--8 & Detailed answer with named evidence, geographic processes, balance, and
clear evaluation. A supported conclusion is present. \\
9--10 & Sustained, precise, and balanced argument with strong case evidence,
clear causal chains, and explicit judgment that answers the command term. \\
\bottomrule
\end{tabularx}
\rowcolors{2}{}{}

\section{Final Error Log}

\begin{tabularx}{\textwidth}{L{0.18\textwidth}Y Y Y}
\toprule
\textbf{Error type} & \textbf{What I wrote or missed} & \textbf{Correct version} & \textbf{Exam-day fix} \\
\midrule
Command term & \rule{0pt}{1.35cm} & & \\
Source use & \rule{0pt}{1.35cm} & & \\
Process chain & \rule{0pt}{1.35cm} & & \\
Case evidence & \rule{0pt}{1.35cm} & & \\
Evaluation & \rule{0pt}{1.35cm} & & \\
\bottomrule
\end{tabularx}

\boxnote{SoftGreen}{Carry forward.}{Before the real paper, carry only the
highest-value fixes: exact source evidence in short answers, named examples in
10-mark answers, and a clear judgment for every evaluate, examine, or to what
extent prompt.}

\end{document}
